\documentclass[11pt]{article}
\usepackage{amsmath,amssymb,amsthm}
\usepackage[margin=1.2in]{geometry}
\newtheorem{theorem}{Theorem}
\newtheorem{lemma}[theorem]{Lemma}
\newtheorem{proposition}[theorem]{Proposition}
\newtheorem{corollary}[theorem]{Corollary}
\newtheorem{remark}[theorem]{Remark}
\title{Erd\H{o}s \#377: the carrier criterion, the carry conservation law,\\
and the exact location of the remaining difficulty}
\author{Samuel Lavery}
\date{Draft --- \today}
\begin{document}
\maketitle

\begin{abstract}
\textbf{\#377 is not proved here.}  What is proved: the balance condition is an
exact statement about a single phasor on the carrier --- $n$ is $p$-balanced iff
$\{n/p^{\,i}\}<\cos\theta_p$ for every $i\ge1$, with $\cos\theta_p=\frac{p+1}{2p}$
--- so that payoff and alignment cost are one angle read two ways, and the cost
of freezing rail $p$ at $n$ is exactly $\beta_p\log n$.  From Kummer's carry
count we obtain an exact conservation law, $\sum_p\kappa_p(n)\log p
=\log\binom{2n}{n}$, and from it an unconditional evaluation valid for
\emph{every} $n$: the $\theta$-weighted mass of balanced primes above
$\sqrt{2n}$ equals $2n(1-\log2)+o(n)$.  We then show this law is
\emph{orthogonal} to \#377 --- its support on small primes is $O(\sqrt n)$ out of
$2n\log2$, measured share $0.0127\to0.00066$ --- so it cannot bound $E$.  Two
candidate reductions are tested and one is refuted: the entropy budget
$\sum_{p\in S(n)}\beta_p\le1$ is \textbf{false} for every fixed truncation
(measured maximum $7.03$), holding only as $d_p\to\infty$.  What survives is a
clean reduction with an explicit constant, and a difficulty localised to a
single range of primes, stated in \S6.
\end{abstract}

\section{The carrier criterion}

For an odd prime $p$ let $D_p=\{0,\dots,\frac{p-1}2\}$, and call $n$
\emph{$p$-balanced} ($n\in G_p$) if every base-$p$ digit of $n$ lies in $D_p$;
by Kummer this is exactly $p\nmid\binom{2n}{n}$.  Write
\[
  \cos\theta_p=\frac{|D_p|}{p}=\frac{p+1}{2p},\qquad \theta_p\nearrow\frac\pi3,
  \qquad \frac\pi3-\theta_p=\frac{1}{\sqrt3\,p}\bigl(1+O(p^{-1})\bigr).
\]

\begin{theorem}[carrier criterion: the universal half-turn]\label{thm:carrier}
Let $F(x)=\mathbf 1[\{x\}<\tfrac12]$.  For every $n\ge1$ and every odd prime
$p$,
\[
  n\in G_p
  \iff
  n\bmod p^{\,j}<\tfrac12p^{\,j}\ \ \forall j\ge1
  \iff
  F\!\Bigl(\frac{n}{p^{\,j}}\Bigr)=1\ \ \forall j\ge1
  \iff
  \Im\, e\!\Bigl(\frac{n}{p^{\,j}}\Bigr)>0\ \ \forall j\ge1 .
\]
\end{theorem}

\begin{proof}
($\Rightarrow$) If every digit is $\le\frac{p-1}2$ then
$n\bmod p^{j}=\sum_{i<j}n_ip^i\le\frac{p-1}2\cdot\frac{p^j-1}{p-1}
=\frac{p^j-1}2<\frac{p^j}2$.
($\Leftarrow$) Induct on $j$.  From $n\bmod p<\frac p2$ we get
$n_0\le\frac{p-1}2$.  Given $n_0,\dots,n_{j-1}\le\frac{p-1}2$, the hypothesis
$n\bmod p^{j+1}\le\frac{p^{j+1}-1}2$ gives
$n_jp^j\le\frac{p^{j+1}-1}2$, so $n_j\le\frac{p^{j+1}-1}{2p^j}<\frac p2$, i.e.\
$n_j\le\frac{p-1}2$.  The last two equivalences are
$\{x\}<\frac12\iff\sin2\pi x>0$.
\end{proof}

\emph{Verified}: $480\,000$ pairs $(n,p)$, $p\in\{3,\dots,1009\}$, zero
mismatches against the digit criterion.

The condition has no digits in it and, more to the point, \emph{no per-rail
constant}: one universal function $F$, threshold exactly the half-turn, applied
at the scales $p^{\,j}$.  The rails differ only in their multiplier.  Since
$F=\frac12+\frac2\pi\sum_{k\ \mathrm{odd}}\frac{\sin2\pi kx}{k}$ carries odd
harmonics only, each with $1/k$ decay, the second (even, non-decaying) lane
seen in the residue-wise reading is the aliasing of that tail onto
$\mathbb Z/p$, not a feature of the carrier --- though it is \emph{not} removed
by the reformulation: the two readings give literally the same subset of
$\mathbb Z/p$, hence identical Fourier coefficients, checked to eight places.

Writing $\cos\theta_p=\frac{p+1}{2p}$ for the density of $D_p$
(so $\theta_p\nearrow\frac\pi3$ and $\cos\theta_p-\cos\frac\pi3=\frac1{2p}$
exactly), payoff and cost are one angle:
\[
  \frac1p=\frac{2}{p}\Bigl(\cos\theta_p-\cos\tfrac\pi3\Bigr)\cdot 1,
  \qquad
  \beta_p:=1-\frac{\log|D_p|}{\log p}=\frac{\log\sec\theta_p}{\log p}.
\]

\begin{proposition}[cost identity]\label{prop:cost}
Let $d_p=\lfloor\log n/\log p\rfloor$.  The density of $G_p$ among the residues
mod $p^{\,d_p+1}$ is $\cos^{d_p+1}\theta_p$, so the cost of freezing rail $p$ at
height $n$ is
\[
  -\log\bigl(\text{density}\bigr)=(d_p+1)\log\sec\theta_p=\beta_p\log n+O(\log p).
\]
\end{proposition}

Write $S(n)=\{p\le n:\ n\in G_p\}$ and
$E(n)=\sum_{p\in S(n)}1/p=\operatorname{pay}(S(n))$.

\section{Band truncation (unconditional)}

\begin{theorem}\label{thm:trunc}
For every $n$ and every integer $K\ge1$,
\[
  E(n)\ \le\ \log(K+1)\ +\ \sum_{p\le n^{1/K}}\frac1p\mathbf 1[n\in G_p]
  \ +\ O\!\Bigl(\frac1{\log n}\Bigr).
\]
\end{theorem}

\begin{proof}
Bound band $k=\{p\in(n^{1/(k+1)},n^{1/k}]\}$ for $k<K$ by its full Mertens mass
$\log\frac{k+1}{k}+O(e^{-c\sqrt{\log n}})$; the sum telescopes to $\log K$, and
the band containing $n^{1/K}$ contributes at most $\log\frac{K+1}{K}$.
\end{proof}

No equidistribution enters, and no band is discarded.  The whole of \#377 is
therefore the deep sum $\sum_{p\le n^{1/K}}\frac1p\mathbf1[n\in G_p]$.

\section{What a budget would buy}

\begin{lemma}[payoff/cost monotonicity]\label{lem:ratio}
With $\beta(x)=\log\frac{2x}{x+1}/\log x$, the ratio
$r(x)=\frac{1/x}{\beta(x)}=\frac{\log x}{x\log\frac{2x}{x+1}}$ is strictly
decreasing on $(e,\infty)$ --- in particular on the odd primes --- and
$r(x)\sim\frac{\log x}{x\log2}\to0$.
\end{lemma}

\begin{proof}
Write $r=uv$ with $u(x)=\frac{\log x}{x}$ and $v(x)=1/g(x)$, where
$g(x)=\log\frac{2x}{x+1}=\log2-\log(1+\frac1x)$.  Both are positive for $x>1$,
since $1+\frac1x<2$ gives $g>0$.  Now $u'(x)=(1-\log x)/x^2<0$ for $x>e$, so
$u$ strictly decreases there; and $x\mapsto\log(1+\frac1x)$ strictly decreases,
so $g$ strictly increases and $v=1/g$ strictly decreases.  A product of two
positive strictly decreasing functions strictly decreases.  The asymptotic
follows from $g(x)\to\log2$.
\end{proof}

\begin{theorem}[LP reduction]\label{thm:lp}
Let $\Lambda(B)=\sup\{\operatorname{pay}(S):\ S\ \text{a finite set of odd primes},\
\sum_{p\in S}\beta_p\le B\}$.  If for some $K$ and some $B<\infty$
\[
  \sum_{p\in S(n),\ p\le n^{1/K}}\beta_p\ \le\ B
  \qquad\text{for all sufficiently large }n,
  \tag{$\mathrm B_{K,B}$}
\]
then $\limsup_n E(n)\le\log(K+1)+\Lambda(B)$; in particular \#377 holds.
Because $\frac{1/p}{\beta_p}=\frac{\log p}{p\log\frac{2p}{p+1}}$ is strictly
decreasing, $\Lambda$ is computed greedily; the integral values are
$\Lambda(1)=0.676190$ (rails $3,5,7$), $\Lambda(2.5)=1.033439$ (rails up to
$29$), $\Lambda(7)=1.417404$ (rails up to $173$), against the fractional
relaxation $\Lambda^{\mathrm{LP}}(1)=0.685523$.
\end{theorem}

\begin{proof}
The bound $\limsup E\le\log(K+1)+\Lambda(B)$ is immediate from
Theorem~\ref{thm:trunc} and the definition of $\Lambda$.  For the greedy
evaluation, consider the fractional relaxation
\[
  \Lambda^{\mathrm{LP}}(B)=\max\Bigl\{\sum_p\frac{x_p}{p}\ :\
  \sum_px_p\beta_p\le B,\ x_p\in[0,1]\Bigr\}\ \ge\ \Lambda(B),
\]
and write $r(p)=\frac{1/p}{\beta_p}$, strictly decreasing by
Lemma~\ref{lem:ratio}.  Let $x$ be feasible and suppose it is not greedy: there
are primes $p<q$ with $x_p<1$ and $x_q>0$.  For small $\delta>0$ set
$x_q'=x_q-\delta$ and $x_p'=x_p+\delta\beta_q/\beta_p$, leaving all other
coordinates fixed.  This is again feasible (the budget used is unchanged, and
$x_p'\le1$, $x_q'\ge0$ for $\delta$ small), and the objective changes by
\[
  \frac{\delta\beta_q}{\beta_p}\cdot\frac1p-\frac{\delta}{q}
  \;=\;\delta\beta_q\Bigl(\frac{1/p}{\beta_p}-\frac{1/q}{\beta_q}\Bigr)
  \;=\;\delta\beta_q\bigl(r(p)-r(q)\bigr)\;>\;0 ,
\]
strictly, since $r$ is strictly decreasing and $\beta_q>0$.  Hence no
non-greedy point is optimal, so the optimum fills the coordinates in increasing
order of $p$ until the budget is exhausted --- which is the greedy rule.  The
integral values follow by evaluating the greedy solution and rounding down at
the last prime.
\end{proof}

\begin{theorem}[$\mathrm B_{K,1}$ is false]\label{thm:budgetfalse}
For every fixed $K$, $(\mathrm B_{K,1})$ fails.  Over the dyadic window
$[N/2,N]$ the maximum of $\sum_{p\le z,\ n\in G_p}\beta_p$ is
\[
\begin{array}{r|cccc}
 & z=7 & z=31 & z=101 & z=1009\\\hline
 N=10^6    & 0.686 & 1.261 & 1.879 & 6.803\\
 N=10^7    & 0.657 & 1.133 & 1.735 & 7.031\\
 N=4\times10^7 & 0.686 & 1.283 & 1.637 & 6.575
\end{array}
\]
and the $90$th percentile at $z=1009$ is already $4.06$--$5.05$.
\end{theorem}

The reason is structural and worth stating, because it says where the budget
\emph{does} live.  The cost identity (Proposition~\ref{prop:cost}) charges
$\beta_p\log n$ per rail only when the rail is asked to freeze at every one of
its $d_p+1$ scales.  For $p$ near $n^{1/K}$ with $K$ fixed, $d_p=K$ is bounded,
the freeze is cheap, and the number of such rails is large: their costs add
without any of them being individually improbable.  The budget binds only as
$d_p\to\infty$, i.e.\ for $p\le n^{o(1)}$.  Correspondingly the exceptional set
at fixed $z$ is \emph{finite}: at $z=31$ the count of $n$ in the window with
cost $>1$ is $44,\ 39,\ 28$ while the window grows by a factor $40$.

\section{Band one is an exact union of intervals}

Above $\sqrt n$ the balanced condition is not merely an interval union in principle: the
endpoints are explicit rational functions of $n$ and $m=\lfloor n/p\rfloor$.

\begin{theorem}[exact interval criterion]\label{thm:interval}
Let $p>\sqrt n$ be prime and put $m=\lfloor n/p\rfloor$.  Then
\[
  n\in G_p
  \qquad\Longleftrightarrow\qquad
  p\,(2m+1)\ \ge\ 2n+1 \quad\text{and}\quad p\ \ge\ 2m+1 .
\]
\end{theorem}

\begin{proof}
Since $p>\sqrt n$ we have $p^{2}>n$, so $n$ has exactly the two base-$p$ digits
$n=mp+r$ with $r=n-mp$ and $0\le m<p$.  By Kummer, $n\in G_p$ iff both digits are at most
$\tfrac{p-1}{2}$.  The bottom digit condition $r\le\frac{p-1}{2}$ reads $2(n-mp)\le p-1$,
i.e.\ $2n+1\le p(2m+1)$.  The top digit condition $m\le\frac{p-1}{2}$ reads $p\ge2m+1$.
\end{proof}

\emph{Verified}: $1\,134\,661$ pairs $(n,p)$ with $n<4000$ and $p>\sqrt n$ --- no mismatch.

\begin{corollary}[the band-one mass, exactly]\label{cor:band1exact}
For each $m\ge1$ the balanced primes $p$ with $\lfloor n/p\rfloor=m$ are exactly the primes
of
\[
  I_m(n)\;=\;\Bigl[\ \max\Bigl(\tfrac{2n+1}{2m+1},\,2m+1\Bigr),\ \tfrac nm\ \Bigr],
  \qquad\text{so}\qquad
  \sum_{\substack{\sqrt n<p\le n\\ n\in G_p}}\frac1p\;=\;\sum_{m\ge1}\sum_{p\in I_m(n)}\frac1p .
\]
\end{corollary}

\begin{proof}
Immediate from Theorem~\ref{thm:interval}: $\lfloor n/p\rfloor=m$ is equivalent to
$\frac{n}{m+1}<p\le\frac nm$, and the first condition cuts that range at
$\frac{2n+1}{2m+1}$, which lies inside it.
\end{proof}

\emph{Verified}: the interval formula reproduces the direct sum to machine precision at
$n=500,1000,2000,4000$, giving $0.273946$, $0.291506$, $0.337442$, $0.301129$.

\begin{remark}
The endpoints $\frac nm$ and $\frac{2n+1}{2m+1}$ are consecutive mediants, so the balanced
fraction of the $m$-th interval is
$\log\frac{\log(n/m)}{\log((2n+1)/(2m+1))}$, which tends to $\tfrac12\log\frac{m+1}{m}$ for
fixed $m$ as $n\to\infty$: exactly half the band, because the mediant bisects the interval
logarithmically.  The top-digit condition $p\ge2m+1$ binds only for $m\gtrsim\sqrt{n/2}$,
i.e.\ only at the very bottom of the band.
\end{remark}

\section{The carry conservation law}

Kummer's criterion in its quantitative form is $v_p\binom{2n}{n}=\kappa_p(n)$,
the number of carries in the base-$p$ addition $n+n$; equivalently
$\kappa_p(n)=\frac{2s_p(n)-s_p(2n)}{p-1}\ge0$, and $n\in G_p\iff\kappa_p(n)=0$.
Unlike the indicator, $\kappa_p$ is \emph{one-signed and additive}, and it obeys
an exact global identity.

\begin{theorem}[conservation]\label{thm:cons}
For every $n\ge1$,
\[
  \sum_{p\le2n}\kappa_p(n)\log p=\log\binom{2n}{n}.
\]
Consequently, for every $n$,
\[
  \sum_{\substack{\sqrt{2n}<p\le2n\\ p\,\nmid\,\binom{2n}{n}}}\log p
  \;=\;2n\,(1-\log 2)\;+\;O\!\bigl(\sqrt n\,\log n\bigr)
  \;=\;0.613706\ldots\,n\;+\;o(n).
\]
\end{theorem}

\begin{proof}
The identity is Legendre's formula summed over $p$.  For $p>\sqrt{2n}$ the
addition has at most two digits, so $\kappa_p\in\{0,1\}$; for $p\le\sqrt{2n}$,
$\kappa_p\le\log(2n)/\log p$, whence
$\sum_{p\le\sqrt{2n}}\kappa_p\log p\le\psi(\sqrt{2n})=O(\sqrt n)$.  Therefore
$\sum_{\sqrt{2n}<p\le2n}\mathbf1[\kappa_p\ge1]\log p=\log\binom{2n}{n}+O(\sqrt n)
=2n\log2+O(\sqrt n\log n)$, and subtracting from
$\theta(2n)-\theta(\sqrt{2n})=2n+o(n)$ gives the display.
\end{proof}

\emph{Measured}, ratio of the left side to $2n(1-\log2)$ at
$n=10^4,10^5,10^6,3\times10^6$: $0.9771,\ 0.9939,\ 0.9979,\ 0.9985$.

This is a genuine common-mode evaluation: exact, uniform in $n$, and free of any
equidistribution input.  It also settles the question it raises.

\begin{proposition}[the law is orthogonal to \#377]\label{prop:orth}
The share of $\log\binom{2n}{n}$ carried by primes $p\le\sqrt{2n}$ is
$O(\sqrt n/n)\to0$ --- measured $0.01268,\ 0.00373,\ 0.00106,\ 0.00066$ at
$n=10^4,10^5,10^6,3\times10^6$.  Hence Theorem~\ref{thm:cons} constrains only
the range $p>\sqrt{2n}$, where by Theorem~\ref{thm:trunc} the Mertens mass is
already at most $\log2$ unconditionally.  It gives \emph{no} information about
$p\le n^{1/K}$, which by Theorem~\ref{thm:trunc} is the whole problem.
\end{proposition}

The law is stated and proved here rather than suppressed, because a
conservation identity of exactly this shape is what the method calls for, and
knowing that it is inert is worth more than not having looked: the $\log p$
weight of the identity and the $1/p$ weight of $E$ are dual, and the identity's
mass sits where $E$'s does not.

\section{Falsification register}

Pre-committed, and all currently \emph{unhit}: (i) an $n$ with
$E_{\le y}(n)>\frac13+\frac15+\frac17$ for $y<3.98\times10^{11}$ and infinitely
many such $n$; (ii) a failure of Theorem~\ref{thm:cons}'s ratio to approach $1$;
(iii) a fixed $z$ at which $\#\{n\le N:\operatorname{cost}_{\le z}(n)>1\}$ grows
with $N$.  Hit and retracted this session: the entropy budget
$(\mathrm B_{K,1})$, refuted by Theorem~\ref{thm:budgetfalse}; and an earlier
maximum of $E$ attained at $n=0$, which is balanced at every rail and is an
artifact of an unrestricted range, not a wall.

\section{Where the difficulty is, exactly}

Let $\gamma_k(n)$ be the fraction of band $k$'s Mertens mass that is balanced,
so $E(n)=\sum_k\gamma_k(n)\log\frac{k+1}{k}+O(1/\log n)$.  Band $k$ imposes
$k+1$ conditions, of which the top one is free for all but a $O(1/\log n)$
fraction of the band (Theorem~\ref{thm:carrier} with $i=k+1$ reads
$n/p^{k+1}<\cos\theta_p$, one cut).  Hence the band common mode is
\[
  \mathbb E\,\gamma_k=\cos^{k}\theta_p+o(1)=2^{-k}+o(1),
  \qquad
  \sum_{k\ge1}2^{-k}\log\frac{k+1}{k}=C_0=0.507834\ldots,
\]
which is the mean of $E$ --- the band decomposition and the global common mode
agree exactly.  \emph{Measured} at $N=10^6$ over $1500$ sampled
$n\in[N/2,N]$:
\[
\begin{array}{r|rrrrr}
 k & \#p & \text{band mass} & \operatorname{mean}\gamma_k & 2^{-k} & \max\gamma_k\\\hline
 1 & 78330 & 0.68925 & 0.47664 & 0.5 & 0.49747\\
 2 & 143 & 0.39526 & 0.22686 & 0.25 & 0.36799\\
 3 & 14 & 0.23712 & 0.13231 & 0.125 & 0.44554\\
 4 & 5 & 0.22167 & 0.06347 & 0.0625 & 0.58912
\end{array}
\]
Band~$1$ is concentrated to $\pm0.02$ over the whole sample; $\max_n\gamma_k$
rises toward $1$ only as the band's prime count falls.  No sampled $n$ raised
more than four of five bands above twice their common mode, and
$\max_n\sum_{k\le6}\gamma_k\log\frac{k+1}{k}=0.6806$ against the trivial
$1.7111$.  \textbf{Boundedness of $E$ is therefore a joint constraint across
near-independent channels, not a per-band decay:} each $\gamma_k$ can
individually approach $1$, and what fails is doing so at many $k$ for one $n$.

The prime count of band $k$ is $\asymp k\,n^{1/k}/\log n$, which is large
exactly when $p\gtrsim\log n$.  This splits \#377 into two halves, and both are
open:
\[
  \boxed{\ \text{(a) } p\le\log n:\ \pi(\log n)\ \text{rails, mass }\log\log\log n\ }
  \qquad
  \boxed{\ \text{(b) } p>\log n:\ \text{concentration of }\gamma_k\ }
\]
Half (a) is where the budget binds ($d_p\ge\log n/\log\log n$), and the ceiling
of the companion note caps it at $\Lambda(1)=0.676190$ --- conditional on
transversality for subsets of $\{p\le\log n\}$, a growing family.  Half (b)
needs $\gamma_k(n)\ll1/(\log k)^{1+\delta}$ uniformly in $n$, which is far
weaker than the truth $2^{-k}$ but still an equidistribution statement: for
fixed $n$, the points $\bigl(\{n/p\},\dots,\{n/p^{\,k}\}\bigr)$ must
equidistribute in $[0,1)^{k}$ as $p$ ranges over the band.  Condition $i$ has
period $t^{\,i+1}/(in)$ in $t$, which exceeds $1$ only for $i\ge k-1$; so at
most two of the $k$ conditions are resolvable at the integer scale, and the rest
are questions about primes in intervals shorter than unit length.  Band~$1$ is
the exception and is proved (companion note, on $m\le n^{(1-\varepsilon)/2}$),
giving $\gamma_1=\frac12+O(\varepsilon)$ --- one band out of infinitely many.

That is the gap, in both halves, named and not papered over.  Nothing in this
note assumes either, and nothing in \S1--\S5 depends on them.

\begin{remark}[register]
\textbf{Proven unconditionally}: Theorems~\ref{thm:carrier},~\ref{thm:interval},
~\ref{thm:trunc}, \ref{thm:lp} (as an implication), \ref{thm:cons};
Corollary~\ref{cor:band1exact}; Propositions~\ref{prop:cost}, \ref{prop:orth}.
Theorem~\ref{thm:interval} and Corollary~\ref{cor:band1exact} give band one exactly: the
balanced primes above $\sqrt n$ are precisely the primes of the explicit intervals
$I_m(n)$, so the ``free half'' of band one is computed rather than estimated.  \textbf{Refuted}: $(\mathrm B_{K,1})$, by measurement.
\textbf{Not proved}: Erd\H{o}s \#377.  \textbf{The literature gate has not been
run}: Erd\H{o}s--Graham--Ruzsa--Straus (1975) and successors must be read at
source before any novelty claim; Theorem~\ref{thm:cons} in particular is an
immediate consequence of Legendre's formula and should be presumed known.
\end{remark}

\end{document}
